%============================================================
\chapter{企業部門横断の実証：法人企業統計}
\label{part:industry}
%============================================================

\section{位置づけ}

第\ref{part:theory}部で定義した $\CCC$ を用いて、業界の状態を判定できるかを検討する。
本章では二つの事前登録済み予測を検証する。結論を先に述べると、
\textbf{両方とも棄却された}。予測Fは符号が逆であり、
予測Eは主張の後半（先行性）が完全に消滅した。

\section{ホットとコールドは一次元ではない}

測定の前に、判定対象を確定する。
式\eqref{eq:gstar}の $g^\star=m/\CCC$ より、成長と資金繰りは独立であった。
同様に、\textbf{業界が熱いことと、そこへ参入して得することは別である}。

需要が逼迫していても、その需要への到達権を仲介者が押さえていれば、
余剰は参入した側に残らない。第\ref{sec:platform-fee}章では
$T_{\rm acq}$ の価格が20ポイントに達することを見る。
したがって二軸で考える必要がある。

\begin{center}
\begin{tabularx}{\textwidth}{@{}llX@{}}
\toprule
軸 & 指標 & 測定可能性 \\
\midrule
需要の逼迫 & $\kappa$ の変化、成長率 & 法人企業統計で測定可能 \\
余剰の取りやすさ & $T_{\rm acq}$ の価格、$\phi_{\rm barg}$ の集中度 & 業種横断の指標が存在しない \\
\bottomrule
\end{tabularx}
\end{center}

本章で扱えるのは前者のみである。四象限のうち上下は区別できるが左右は区別できない。

\begin{remark}[横軸が測れない構造的な理由]
\label{rem:no-horizontal}
これは資料の不足ではなく、分類の非対称による。

業種分類は参入者側の目線で切られている。ある事業者が情報通信業であれば
情報通信業として集計される。ところが\textbf{仲介者は業種を跨ぐ}。
決済基盤は情報通信業に属しながら、小売・飲食・旅行・教育のすべてに対して仲介する。

したがって $\phi_{\rm barg}^{\text{仲介}}$ を業種別に足し上げても、
仲介者の取り分はどこか一つの業種に計上されるだけであり、
\textbf{どの業種から取ったかは復元できない}。
第\ref{sec:three-ops}節の記法でいえば、仲介者は主体の添字 $i$ の一つであって、
$\Phi$ の類型でも業種でもない。

参入者側の利益率には仲介への流出が既に反映されているが、
それが流出によるのか他の要因によるのかは業種別集計から分離できない。
分離には仲介者側から見た断面が必要であるが、
それが得られない理由は第\ref{sec:pred-J}節で述べる。
\end{remark}


\begin{center}
\begin{tabularx}{\textwidth}{@{}llX@{}}
\toprule
\multicolumn{3}{@{}l}{\textbf{本章で扱う理論量}}\\
\midrule
$\kappa_C$ & 式\eqref{eq:kappa} & DSO。35--71 日（規模別） \\
$\kappa_S$ & 式\eqref{eq:kappa} & $-$DPO。$-21$--$-44$ 日 \\
$W,\ r$ & 第\ref{sec:growth}章 & 運転資本、売上高 \\
$\CCC$ & 式\eqref{eq:ccc} & 全産業 37 日、業種別 5--210 日 \\
$g^\star$ & 式\eqref{eq:gstar} & $m$ 未取得のため仮定値で例示 \\
$m$ & 第\ref{sec:growth}章 & \textbf{未取得} \\
\bottomrule
\end{tabularx}
\end{center}

\section{データ源}

法人企業統計調査の時系列データ（e-Stat 統計表ID 0003060791）を用いた。
統計法に基づく基幹統計調査であり、年次別調査は全ての営利法人等を対象とする。
資本金階層は2百万円未満まで分かれており、
第\ref{sec:frame}節で問題とした上場企業偏重が生じない。

取得したのは受取手形、売掛金、製品又は商品、支払手形、買掛金、売上高の6項目である。
既製の回転期間指標ではなく原数値を取得したのは、
統計側の定義（分母が売上高か売上原価か）に依存せず
第\ref{sec:growth}章の定義をそのまま適用するためである。

\begin{align}
\mathrm{DSO} &= \frac{\text{受取手形}+\text{売掛金}}{\text{売上高}}\times 365 \\
\mathrm{DPO} &= \frac{\text{支払手形}+\text{買掛金}}{\text{売上高}}\times 365 \\
\CCC &= \mathrm{DSO} + \mathrm{DIO} - \mathrm{DPO}
\end{align}

これらは式\eqref{eq:kappa}の信用ポジションを日数で表したものである。すなわち
\begin{equation}
\kappa_C \;\leftrightarrow\; \mathrm{DSO},\qquad
\kappa_S \;\leftrightarrow\; -\mathrm{DPO},\qquad
\sum_i \kappa_i \;\leftrightarrow\; \mathrm{DSO}-\mathrm{DPO}
\end{equation}
と対応する。$\sum_i\kappa_i>0$ であれば与信の出し手である。
また式\eqref{eq:ccc}の運転資本は $W = (\text{DSO}+\text{DIO}-\text{DPO})\times r/365$ にあたる。

\section{予測Fの検証：棄却}

\subsection{予測}

\begin{quote}
\textbf{予測F}：小規模階層ほど $\mathrm{DSO}-\mathrm{DPO}$ が大きい。
すなわち小規模ほど与信の出し手になっている。
\end{quote}

根拠は第\ref{sec:layers}章の「信用層は交渉力の写像である」という主張であった。
交渉力のない小規模企業は支払いを待たされる側に置かれるはずだ、という推論である。

\subsection{結果}

\begin{center}
\begin{tabular}{@{}lrrrr@{}}
\toprule
規模（資本金） & DSO & DIO & DPO & DSO$-$DPO \\
\midrule
10億円以上 & 71 & 20 & 44 & \textbf{27} \\
1億円以上 -- 10億円未満 & 61 & 20 & 47 & 15 \\
5千万円以上 -- 1億円未満 & 47 & 20 & 37 & 10 \\
2千万円以上 -- 5千万円未満 & 49 & 15 & 39 & 10 \\
1千万円以上 -- 2千万円未満 & 53 & 18 & 35 & 18 \\
1千万円未満 & 35 & 12 & 21 & \textbf{14} \\
\bottomrule
\end{tabular}
\end{center}
\begin{center}\footnotesize 全産業（除く金融保険業）、2024年度、単位は日\end{center}

\begin{figure}[htbp]
\centering
\insertfig{ccc-trend}
\caption{資本金階層別の DSO と DPO（全産業、2024年度）。小規模階層は双方とも短い。}
\label{fig:size-dso-dpo}
\end{figure}

\textbf{予測Fは棄却される。符号が逆である。}
10億円以上が27日、1千万円未満が14日であり、大企業のほうが与信の出し手である。
この順序は2015年から2024年まで10年間一貫しており、
大企業側はむしろ19日から27日へ拡大している。

\subsection{推論の誤り}

誤りは特定できる。\textbf{DSO と DPO が両方とも動くことを見落としていた。}

小規模企業は売掛金を待たされる（DSO が不利になる）。
しかし同時に、仕入先からも信用されないため買掛金を伸ばせない。
図\ref{fig:size-dso-dpo}のとおり、1千万円未満の DPO は21日であり、
10億円以上の44日の半分以下である。

交渉力の欠如は両側に作用する。回収も遅いが支払いも早い。
差し引きすると\textbf{交渉力の弱さは正味の与信ポジションには現れない}。

\begin{proposition}[与信を出すには資本が要る]
第\ref{part:solo}章で $E$ と $\kappa<0$ が代替物であることを示した。
逆向きも成立する。$E$ があれば $\kappa>0$ を保有できる。
DSO 71日という長期の回収を許容できるのは、それを支える資本があるためである。
\end{proposition}

\subsection{業種によって符号が割れる}

\begin{center}
\begin{tabular}{@{}lrrr@{}}
\toprule
業種 & 10億円以上 & 1千万円未満 & 差 \\
\midrule
情報通信業 & 64 & 15 & $+49$ \\
建設業 & 67 & 19 & $+48$ \\
全産業 & 27 & 14 & $+13$ \\
製造業 & 26 & 24 & $+2$ \\
職業紹介・労働者派遣業 & 24 & 25 & $-1$ \\
小売業 & $-16$ & 6 & $\mathbf{-22}$ \\
\bottomrule
\end{tabular}
\end{center}
\begin{center}\footnotesize DSO$-$DPO、2024年度、単位は日\end{center}

\textbf{小売業のみ符号が反転する。}
大手小売の $-16$ 日は第\ref{part:theory}部で述べたマイナスの信用ポジションそのものである。
即時に現金を回収し、仕入は後払いとする構造であり、
ここでは大企業が与信の受け手、小規模小売が出し手となっている。

情報通信業と建設業の $+49$ 日は、大企業が受託の立場で長期の立替を担うことを示す。
第\ref{part:catalog}部の族1-3（後払単発）および族5-3（コストプラス）の領域である。

したがって「大企業ほど与信の出し手」という命題自体も業種を跨いだ平均にすぎない。
決定しているのは規模ではなく、\textbf{その業種で支配的な $\Phi$} である。
第\ref{part:results}部で「$\kappa$ の水準は業種間で構造的に異なるため横断比較は無意味」と
記録した教訓が、規模軸でも成立していた。

\subsection{第\texorpdfstring{\ref{part:solo}}{VII}部への含意}

第\ref{part:solo}章では、無資産の一人事業には $\kappa_C<0$ が要求されると演繹した。
データはその要求が満たされていないことを示す。
1千万円未満の階層は DSO 35日、DPO 21日で、正味14日の出し手である。

ただし予測Fを立てた際に想定した形とは異なる。
小規模層が極端に $\kappa>0$ なのではなく、\textbf{DSO も DPO も双方が短い}。
取引そのものが小さく短い。実現可能領域の外にいるというより、
別の座標にいるという見え方になる。

\section{予測Eの検証：棄却}
\label{sec:e-verification}

\subsection{予測}

\begin{quote}
\textbf{予測E}：$\CCC$ の変化率は売上高変化率に対して\emph{先行または同期}する。
特に売上債権回転期間の悪化は売上高の減速に\emph{先行}する。
\end{quote}

$\CCC$ の年次変化 $\Delta\CCC(t)$ と売上高成長率 $g(t+k)$ の相関を
業種別に算出し、集約系列と短期系列を除く45業種で集計した。

\subsection{結果}

\begin{center}
\begin{tabular}{@{}lrrrl@{}}
\toprule
ラグ & 平均相関 & 中央値 & 負の割合 & 符号検定 \\
\midrule
先行 $k=-1$ & 0.036 & 0.058 & 47\% & $p=0.72$ \\
同期 $k=0$ & $-0.159$ & $-0.228$ & 69\% & $\mathbf{p=0.008}$ \\
遅行 $k=+1$ & 0.072 & 0.118 & 31\% & --- \\
遅行 $k=+2$ & 0.085 & 0.035 & 44\% & --- \\
\bottomrule
\end{tabular}
\end{center}

\textbf{同期は支持された。} $\CCC$ の悪化と売上高の減速は同一年度に生じる。
45業種中31業種で負の相関、符号検定で $p=0.008$ である。

\textbf{先行は完全に棄却された。} 47\% が負というのは無情報である（$p=0.72$）。
$\CCC$ の変化は翌年度の売上高について何の情報も持たない。

\subsection{同期の負相関は分母効果かもしれない}

機構が想定と逆であった可能性が高い。
$\CCC = W/r$ は分母に売上高を持つ。
売上が減少すれば、在庫と売掛金が同速度で減少しない限り $\CCC$ は\emph{定義上}悪化する。

第\ref{part:results}部で族2の媒体別パネルを扱った際、
これを対抗仮説として明示した。発行額が縮小して比率が機械的に上昇する、という論点である。
本章では第\ref{sec:denominator-effect}節でこの分離を行い、
\textbf{分母効果の寄与は17\%程度}であることを示す。

\begin{center}
\begin{tabularx}{\textwidth}{@{}Xl@{}}
\toprule
主張 & 判定 \\
\midrule
$\CCC$ の変化が売上高変化に同期する & 支持。分母効果の寄与は17\%（第\ref{sec:denominator-effect}節） \\
$\CCC$ の変化が売上高変化に先行する & \textbf{棄却}（$p=0.72$） \\
$\CCC$ が景気先行指標になる & \textbf{棄却} \\
\bottomrule
\end{tabularx}
\end{center}

\subsection{符号が業種で割れることは発見である}

同期相関で $|r|>0.4$ となる業種は45中13である。

\begin{center}
\begin{tabular}{@{}lrlr@{}}
\toprule
負に強い業種 & $r$ & 正に強い業種 & $r$ \\
\midrule
パルプ・紙・紙加工品製造業 & $-0.78$ & 農業、林業 & $+0.80$ \\
金属製品製造業 & $-0.67$ & その他の輸送用機械器具製造業 & $+0.69$ \\
情報通信業 & $-0.58$ & 宿泊業、飲食サービス業 & $+0.54$ \\
生産用機械器具製造業 & $-0.53$ & 農林水産業 & $+0.52$ \\
\bottomrule
\end{tabular}
\end{center}

負に強い業種はいずれも受注生産型で在庫と売掛金が大きく、分母効果が強く出る条件と一致する。

一方、正の業種が存在することが重要である。
\textbf{分母効果のみであれば全業種で負になるはずである。}
正の業種は、売上が伸びるときに $\CCC$ も伸びる構造、
すなわち成長のために運転資本を積む側であり、
式\eqref{eq:gstar}の $g^\star$ が拘束している領域にあたる。

したがって $\CCC$ の変化には分母効果以外の成分も含まれる。
ただし本データでは分離できていない。

\section{予期しなかった観測：\texorpdfstring{$\CCC$}{CCC} の長期上昇}

検証とは別に、事前に予測していなかった事実が観測された。

\begin{figure}[htbp]
\centering
\insertfig{lag-corr}
\caption{全産業（除く金融保険業）の $\CCC$ と正味与信ポジションの推移。}
\label{fig:ccc-trend}
\end{figure}

全産業の $\CCC$ は2000年度の24.9日から2024年度の37.2日へ、
25年間でほぼ単調に増加している。正味与信ポジションも9.8日から18.5日へ倍増した。
\textbf{日本の企業部門全体が、運転資本をより多く抱える方向へ動いている。}

\subsection{ただし成長制約は悪化していない}

この観測から成長制約の悪化を結論したくなる。
式\eqref{eq:gstar}の分母が5割増えているからである。
\textbf{この推論は誤りである。}

同期間に営業利益率 $m$ は 2.62\% から 5.01\% へ 91\% 改善しており、
分母の悪化を上回る。結果として比 $m/\CCC$ は 38.5\% から 49.2\% へ上昇した
（命題\ref{prop:gstar-flat}）。

\begin{center}
\begin{tabular}{@{}lrrr@{}}
\toprule
 & 2000年度 & 2024年度 & 変化 \\
\midrule
$\CCC$ & 24.9 日 & 37.2 日 & $+49\%$ \\
$m$ & 2.62\% & 5.01\% & $+91\%$ \\
$m/\CCC$ & 38.5\% & 49.2\% & $+28\%$ \\
\bottomrule
\end{tabular}
\end{center}

ここで用いたのは $m/\CCC$ であり、
系\ref{cor:gstar}の $g^\star=(m+d-I/r)/\CCC$ ではない
（注\ref{rem:gstar-flat-caveat}）。
$d$ と $I$ を長期に遡って取得していないためである。

\begin{remark}[比を扱う際の規律]
$\CCC$ は比の分母であり、それ単独では成長制約について何も述べない。
本節が当初 $\CCC$ の上昇のみを観測して停止していたのは、
$m$ を取得していなかったためである。
\textbf{比の一方のみを測って結論を導くことは、
第\ref{sec:narrative}節が戒める事後的物語化と同型の誤りである。}
異なるのは、後者が解釈の後付けであるのに対し、
これは測定の不足を推論で埋めようとする点にある。
\end{remark}

$\CCC$ 上昇そのものの要因は依然として未検証である。
産業構成の変化、支払サイト規制、手形の減少、金利水準のいずれが寄与しているかは
本稿では分離していない。ただしこの問いの重要性は、
$g^\star$ が悪化していない以上、当初想定より小さい。

\section{業種別の水準（参考）}

第\ref{part:results}部の教訓どおり水準の横断比較は解釈できないが、
$\Phi$ の分布を示す参考として記録する。

\begin{center}
\begin{tabular}{@{}lrlr@{}}
\toprule
$\CCC$ が長い業種 & 日 & $\CCC$ が短い業種 & 日 \\
\midrule
リース業 & 210 & 飲食サービス業 & 5 \\
不動産業 & 89 & 水運業 & 7 \\
化学工業 & 80 & その他の運輸業 & 9 \\
繊維工業 & 78 & 娯楽業 & 9 \\
はん用機械器具製造業 & 77 & 自動車・同附属品製造業 & 9 \\
生産用機械器具製造業 & 64 & 電気業 & 10 \\
\bottomrule
\end{tabular}
\end{center}
\begin{center}\footnotesize 2024年度、集約系列を除く\end{center}

序列は第\ref{part:catalog}部の類型と整合する。
リース業（族3-1）と不動産業は資産を先に取得して時間で回収する構造であり、
飲食サービス業（族1-1、即時交換）は $\kappa\approx 0$ に張り付く。
第1回に述べた業種差の構造が、そのまま現れている。

\section{\texorpdfstring{$\CCC$}{CCC} 上昇の要因分解}
\label{sec:ccc-decomposition}

前節の $\CCC$ 上昇について、要因を分解する。
事前に次の予測を登録した。

\begin{quote}
\textbf{予測M}：$\CCC$ の上昇は、業種内効果より構成効果の寄与が大きい。
$\CCC$ の長い業種の比重が上昇したためである。

\textbf{対抗仮説}：業種内効果が支配的なら、構成とは無関係に
各業種で運転資本が増えていることになる。
\end{quote}

\subsection{シフト・シェア分解}

全体の $\CCC$ を業種別の売上加重平均として書き、差分を三項に分ける。
\begin{equation}
\Delta\CCC = \underbrace{\sum_k \bar{w}_k\,\Delta\CCC_k}_{\text{業種内効果}}
 + \underbrace{\sum_k \Delta w_k\,\bar{\CCC}_k}_{\text{構成効果}}
\label{eq:shift-share}
\end{equation}

\begin{center}
\begin{tabular}{@{}lrrr@{}}
\toprule
期間 & $\Delta\CCC$ & 業種内効果 & 構成効果 \\
\midrule
2000$\to$2024（共通33業種） & $+11.4$ 日 & $+9.5$ 日（82.9\%） & $+1.9$ 日（17.1\%） \\
2010$\to$2024（共通45業種） & $+10.1$ 日 & $+9.9$ 日（98.1\%） & $+0.2$ 日（1.9\%） \\
\bottomrule
\end{tabular}
\end{center}

\textbf{予測Mは棄却される。} 業種内効果が8割から9割8分を占める。
分類が安定している2010年以降では構成効果はほぼゼロである。
サービス業化による説明は成立しない。

構成効果が小さい理由も明快である。
情報通信業が比重0.2\%から5.8\%へ増えて $+3.05$ 日を押し上げる一方、
電気機械器具製造業が6.3\%から2.1\%へ縮んで $-1.62$ 日、
広告業が5.6\%から0.8\%へ縮んで $-1.59$ 日となり、
\textbf{押し上げと押し下げがほぼ相殺する}。

\subsection{業種内効果の成分分解}

業種内効果を $\mathrm{DSO}$、$\mathrm{DIO}$、$\mathrm{DPO}$ に分ける。

\begin{center}
\begin{tabular}{@{}lrrrl@{}}
\toprule
期間 & DSO & DIO & DPO & 主因 \\
\midrule
2000$\to$2024 & $-0.93$ & $+3.49$ & $\mathbf{+6.90}$ & DPO の短縮 \\
2010$\to$2024 & $\mathbf{+5.26}$ & $+4.02$ & $+0.62$ & DSO の伸長 \\
\bottomrule
\end{tabular}
\end{center}
\begin{center}\footnotesize 単位は $\CCC$ への寄与日数\end{center}

系\ref{cor:ccc-change}に照らすと、DSO と DPO は $\tau_C$ と $\tau_S$、
DIO は $\tau_{\mathrm{inv}}$ に対応する。したがって

\begin{center}
\begin{tabular}{@{}lrr@{}}
\toprule
期間 & 契約の変化 $\sum\Delta\tau_i$ & 運用の変化 $\Delta\tau_{\mathrm{inv}}$ \\
\midrule
2000$\to$2024 & $+5.97$ 日 & $+3.49$ 日 \\
2010$\to$2024 & $+5.88$ 日 & $+4.02$ 日 \\
\bottomrule
\end{tabular}
\end{center}

\textbf{両期間とも契約の変化が6割前後を占める。}
$\CCC$ の上昇は、需要や在庫管理の変化より契約条件の変更に由来する。

\textbf{期間によって主因が入れ替わる。} 手形の推移がこれを説明する。

\begin{center}
\begin{tabular}{@{}lrrr@{}}
\toprule
 & 2000 & 2010 & 2024 \\
\midrule
受取手形 & 12.1 日 & 6.1 日 & 4.7 日 \\
支払手形 & 15.3 日 & 7.2 日 & 5.1 日 \\
\bottomrule
\end{tabular}
\end{center}

2000年代に手形は三分の一以下へ減少した。
支払手形が15.3日から7.2日へ8.1日分消えており、
DPO の短縮6.9日はほぼこれで説明がつく。
\textbf{手形は $\kappa_S<0$ を作る手段であり、その消滅が $\CCC$ を押し上げた。}

2010年代には手形の減少が一巡し、DSO が5.26日伸びた。
内訳は受取手形 $-1.37$ 日、売掛金 $+6.62$ 日である。
手形が減り続けているにもかかわらず、売掛金がそれを上回って伸びている。
業種別では45業種中28業種で $\mathrm{DSO}$ が増加し、
卸売業（59$\to$70日）、小売業（22$\to$28日）、
情報通信業（62$\to$78日）、建設業（65$\to$71日）の寄与が大きい。

\begin{remark}[科目移動と分離できない]
\label{rem:densai}
売掛金の伸長には二つの解釈がある。
実質的な回収期間の長期化と、電子記録債権への移行に伴う科目の移動である。
後者であれば $\CCC$ の上昇は見かけ上のものにすぎない。

分離には電子記録債権の残高が必要だが、
法人企業統計の調査項目に含まれない。
第\ref{sec:obstacle-taxonomy}節の分類では（iii）識別不能に属する。
\end{remark}

\section{同期負相関の分離：分母効果の寄与}
\label{sec:denominator-effect}

第\ref{sec:e-verification}節で観測した同期負相関について、
需要の情報か $\CCC=W/r$ の定義に由来する分母効果かを分離する。

対数差分をとると $\Delta\ln\CCC = \Delta\ln W - \Delta\ln r$ である。
分母効果のみであれば $W$ は動かない。したがって弾力性
\begin{equation}
\beta = \frac{\Delta \ln W}{\Delta \ln r}
\label{eq:wc-elasticity}
\end{equation}
を推定すれば、$1-\beta$ が分母効果の寄与となる。

\begin{quote}
\textbf{予測N}：$\Delta\ln W$ と $\Delta\ln r$ は正の相関を持つ。
すなわち売上が減れば運転資本も減る。
\end{quote}

\subsection{結果}

\begin{center}
\begin{tabularx}{\textwidth}{@{}Xr@{}}
\toprule
プール推定（44業種、946観測） & $\beta = 0.830$ \\
決定係数 & $R^2 = 0.076$ \\
業種別 $\beta$ の中央値 & $0.625$ \\
業種別 $\beta$ の標準偏差 & $1.161$ \\
\midrule
分母効果の寄与 $1-\beta$ & $0.170$ \\
\bottomrule
\end{tabularx}
\end{center}

\textbf{予測Nは支持される。} 運転資本は売上にほぼ比例して動き、
分母効果に帰せられるのは17\%程度である。
第\ref{sec:e-verification}節の同期負相関は、大部分が需要側の情報を含む。

\subsection{遅れ構造の検証：予測OとP}
\label{sec:lag-test}

命題\ref{prop:kappa-tau}によれば $W(t)$ は $r$ の過去 $\tau$ 期間の積分に対応する。
$\CCC$ が長い業種では前年の売上高にも依存するはずである。
次を事前に登録した。

\begin{quote}
\textbf{予測O}：$\Delta\ln r(t-1)$ の係数 $\beta_1$ は、
その業種の $\CCC$ が長いほど大きい。

\textbf{予測P}：ラグを含めた総弾力性 $\beta_0+\beta_1$ は
同時点のみの推定値 $0.830$ より大きく、分母効果の寄与は17\%より小さい。
\end{quote}

分布ラグ模型 $\Delta\ln W(t)=\beta_0\Delta\ln r(t)+\beta_1\Delta\ln r(t-1)+\varepsilon$
を43業種886観測で推定した。

\begin{center}
\begin{tabular}{@{}lr@{}}
\toprule
$\beta_0$（同時点） & $0.836$ \\
$\beta_1$（1期ラグ） & $-0.015$ \\
総弾力性 & $0.821$ \\
$R^2$ & $0.081$ \\
\bottomrule
\end{tabular}
\end{center}

\textbf{予測OとPはいずれも棄却される。}
$\beta_1$ はほぼゼロであり、$\CCC$ との相関は $-0.101$ で符号も逆である。
$\CCC<30$日の業種では $\beta_1=+0.217$、$\CCC\geq 60$日の業種では $-0.267$ であった。
総弾力性は $0.830$ から $0.821$ へ微減し、分母効果の寄与は17\%から18\%へ実質的に不変である。

\begin{remark}[命題\ref{prop:kappa-tau}は棄却されない]
\label{rem:lag-undetectable}
遅れ構造が検出されないことは、式\eqref{eq:kappa-tau}の誤りを意味しない。
$\CCC$ の水準は最大でも210日、大半は30--60日であり、
年次データの粒度では1年内に収まる。
$\beta_1$ が有意となるには $\CCC$ が1年を超える必要があるが、そのような業種は存在しない。

四半期データであれば検出できる可能性があるが、
四半期別調査は標本法人の仮決算に基づくため、決算整理の影響を受ける。
\end{remark}

\begin{remark}[確度は限定的である]
$R^2=0.076$ は極めて低い。$\Delta\ln W$ の分散の92\%は $\Delta\ln r$ で説明されない。
業種別のばらつきも大きく、$\beta<0.5$ が14業種、$\beta>1.0$ が13業種、
符号が負の業種も存在する（小売業 $-0.43$、リース業 $-0.01$）。

年次の $\Delta\ln W$ には決算期の在庫調整、大口取引の期ずれ、
会計方針の変更が混入する。
\textbf{「分母効果が支配的でない」とは言えるが、
「需要の情報である」と積極的に主張するには不足である。}
\end{remark}

\section{容量制約の簡略化が合わない観測}
\label{sec:cap-observation}

注\ref{rem:cap-simplification}で、$\mathrm{Cap}$ を定数とする簡略化を導入した。
本節はこれが現実と合わない事例を記録する。理論の修正ではなく、限界の記録である。

\subsection{観測}

総合型フィットネスクラブ大手一社の2026年時点の料金体系を確認した。
同社は60歳以上を対象とする平日昼間限定プランを、通常プランと並行して提供している。
両者は週1（月4回まで）、週2（月8回まで）、無制限の三段階で構成が同一であり、
価格のみが異なる。

\begin{center}
\begin{tabular}{@{}lrrr@{}}
\toprule
施設カテゴリ & 週1 & 週2 & 無制限 \\
\midrule
I & 4.0\% & 10.3\% & 22.1\% \\
II & 4.8\% & 9.9\% & 21.5\% \\
III & 4.1\% & 9.6\% & 21.7\% \\
IV & 4.6\% & 10.0\% & 21.5\% \\
\bottomrule
\end{tabular}
\end{center}
\begin{center}\footnotesize 60歳以上プランの通常プランに対する割引率\end{center}

\textbf{四カテゴリすべてで、割引率が利用量とともに単調に増加する。}
週1では約4\%にとどまるが、無制限プランでは約22\%に達する。
週1から無制限への価格倍率は、通常プランで1.67--1.74倍、
60歳以上プランで1.36--1.43倍である。

\subsection{簡略化との不整合}

式\eqref{eq:cap-simplified}の下では、無制限プランを大幅に値引きすることは
容量制約を圧迫するため説明が難しい。
一方、式\eqref{eq:cap-general}の一般形では説明がつく。
平日昼間という時間帯の限定により $\operatorname{supp}\delta_j$ が
$\mathrm{Cap}(t)$ の遊休部分に収まるため、
ピーク時刻の制約を悪化させずに総量を増やせる。

すなわち\textbf{事業者は総量ではなく時間分布を操作している}。
本稿の簡略化はこの手段を表現できない。

\subsection{この観測から理論を修正しない理由}

\begin{remark}[事例からの一般化を行わない]
\label{rem:no-generalization}
本節の観測は一社の料金表に基づく。第\ref{part:theory}部は
全ての $\Phi$ を覆うことを意図した抽象であり、
第\ref{part:results}部は特定の業種と企業に限定された具体である。
\textbf{具体は網羅性を持たない。}

$n=1$ の観測から式\eqref{eq:cap-simplified}を書き換えれば、
一つの事例が27類型すべてに波及する。
これは第\ref{sec:narrative}節が戒める事後的物語化と同型であり、
標本一件から母集団の構造を改訂することにあたる。

理論の修正が正当化されるのは、演繹的な誤りが見つかった場合、
複数の独立な領域で同じ限界が現れた場合、
または第\ref{part:catalog}部の演繹に影響する場合である。
本節はいずれにも該当しない。
したがって観測として記録し、修正は行わない。
\end{remark}

\subsection{予測Lの検証：未達}

第\ref{sec:retiree}節に関連して、次の予測を事前に登録した。

\begin{quote}
\textbf{予測L}：高齢層ほど定額契約より都度払いを選ぶ。
社会情動的選択性理論が予測する「限られた時間展望のもとでは確実なものを選ぶ」
という選好の帰結。定額契約はオプションの束であり、都度払いは確実である。

\textbf{対抗仮説}：高齢層は可処分時間が多く利用回数が増えるため、
定額のほうが有利になる。この説では定額選好が\emph{高まる}。
\end{quote}

符号が逆であるため識別可能な形をしている。

しかし\textbf{消費者側の選択データが得られず、検証できていない}。
必要なのは年齢別の契約形態選択分布であり、事業者は保有するが公開しない。

得られたのは供給側の価格設計のみである。
上記の割引構造は、高齢層がより多く利用するという前提に立っており、
対抗仮説と整合し予測Lとは逆を向く。
ただし\textbf{これは事業者の信念であって消費者の選択ではない}。
企業が需要を正しく読んでいるという前提を要するため、間接証拠にとどまる。

\section{本章の限界}

\begin{enumerate}[label=(\arabic*)]
\item \textbf{分母効果の分離は確度が低い。} 弾力性推定により寄与は17\%程度と得たが、
$R^2=0.076$ であり、$\Delta\ln W$ の変動の大半が未説明である。

\item \textbf{四象限の左右が測れていない。} 余剰の取りやすさについて
業種横断の指標が存在せず、需要側のみの一次元判定にとどまる。

\item \textbf{売掛金伸長の解釈が確定しない。} 第\ref{sec:ccc-decomposition}節により
産業構成は排除し、手形の消滅で半分を説明したが、
2010年代の売掛金伸長が実質的な長期化か科目移動かを分離できていない（注\ref{rem:densai}）。

\item \textbf{予測Lを検証していない。} 第\ref{sec:cap-observation}節のとおり、
消費者側の選択データが得られていない。供給側の価格設計は対抗仮説と整合するが、
間接証拠にとどまる。

\item \textbf{棚卸資産が「製品又は商品」のみである。} 仕掛品と原材料を含めていないため、
DIO は過小である。$\CCC$ の水準は実際より短く出ている。
差分と相関を扱う本章の分析には影響が小さいが、水準の絶対値は解釈できない。
\end{enumerate}

\section{二つの予測が両方棄却されたことについて}

予測Fは符号が逆、予測Eは先行性が消滅した。登録した予測は両方とも外れている。

第\ref{part:results}部で事前登録の効果を記録したが、本章ではそれがより明瞭に働いた。
予測を先に固定していなければ、同期の負相関を
「$\CCC$ は景気指標として機能する」と読み、
大企業の与信ポジションを「大企業が取引先を支えている」と解釈していた可能性が高い。
いずれも後付けの物語である。

なお、予測が外れたことは枠組みの棄却ではない。
$\CCC$、$\kappa$、$\Phi$ という記述は有効に機能しており、
外れたのは\textbf{これらの量が何を予測するかについての推測}である。
記述の道具としての有効性と、予測の道具としての有効性は別である。
